% !PRINTED 445
\weberpage{445}
%<<CONT
into the equations by which the unchanged point is determined [\S.\,116, (3)], we obtain:

\[
\begin{array}{l}
(\beta+1)\,x_1 + \gamma\,\varepsilon^3 x_2 + \alpha\,\varepsilon^2 x_3 = 0,\\
\gamma\,\varepsilon^{-3} x_1 + (\alpha+1)\,x_2 + \beta\,\varepsilon^{-1} x_3 = 0,\\
\alpha\,\varepsilon^{-2} x_1 + \beta\,\varepsilon\, x_2 + (\gamma+1)\,x_3 = 0.
\end{array}
\]
These three equations all reduce to the single one which one derives, e.g., from the first by multiplication by $\beta$ and application of \S.\,115, (8):

\weq{5}{\alpha\,\gamma\,x_1 + \beta\,\gamma\,\varepsilon^3 x_2 + \alpha\,\beta\,\varepsilon^2 x_3 = 0,}
and which represents a straight line. This line is therefore a \gesperrt{principal axis}.

We obtain 21 such principal axes, since we have 21 substitutions of degree 2.

The functions which, set equal to zero, represent the 21 principal axes can easily be formed from (5), if one transforms the function

\weq{6}{A = \alpha\,\gamma\,x_1 + \beta\,\gamma\,\varepsilon^3 x_2 + \alpha\,\beta\,\varepsilon^2 x_3}
by the substitutions $\tau^r$, $\tau^r\chi$, $\tau^r\chi^2$ ($r = 0, 1, 2, 3, 4, 5, 6$). One thus obtains the 21 functions:

\weq{7}{\begin{array}{l}
A_{1,r} = \alpha\,\gamma\,\varepsilon^r x_1 + \gamma\,\beta\,\varepsilon^{2r+3} x_2 + \alpha\,\beta\,\varepsilon^{4r+2} x_3\\
A_{2,r} = \beta\,\gamma\,\varepsilon^{r+3} x_1 + \alpha\,\beta\,\varepsilon^{2r+2} x_2 + \alpha\,\gamma\,\varepsilon^{4r} x_3\\
A_{3,r} = \alpha\,\beta\,\varepsilon^{r+2} x_1 + \alpha\,\gamma\,\varepsilon^{2r} x_2 + \beta\,\gamma\,\varepsilon^{4r+3} x_3,
\end{array}}
from which it is easy to see that the 21 principal axes are in fact all different from one another.

We summarise this as follows:

\noindent 2.\ \gesperrt{There belong 21 different principal axes to the group $G_{168}$, of which each can be transformed into every other by substitutions of the group.}

Since there are only 21 principal axes, each of them must remain unchanged under eight substitutions, and these eight substitutions form a group. To determine this group, we apply to the expression $A$ represented in (6) the substitution $\omega$ [\S.\,115, (5)], and thereby it passes over into

\weq{8}{\begin{array}{l}
\gamma\,(\alpha^2 + \beta^2\varepsilon^3 + \alpha\,\beta\,\varepsilon^2)\,x_1 + \beta\,(\alpha\,\gamma + \gamma^2\varepsilon^3 + \alpha^2\varepsilon^2)\,x_2\\
{}+ \alpha\,(\gamma^2 + \beta\,\gamma\,\varepsilon^3 + \beta^2\varepsilon^2)\,x_3.
\end{array}}
